\chapter{§1.3 运算律}
\label{sec:1.3}


\prerequisite{§1.2 四则运算}

\gradelevel{3-4 年级}


\section*{加法交换律与结合律}


\begin{explore}


小明从家走到学校要经过一个公园。从家到公园是 $300$ 米，从公园到学校是 $500$ 米。从家到学校一共走多少米？如果反过来，从学校走到家呢？

\end{explore}

\begin{understand}


\textbf{加法交换律}：两个数相加，交换加数的位置，和不变。

\[
a + b = b + a
\]


例如： $300 + 500 = 500 + 300 = 800$ 。

\textbf{加法结合律}：三个数相加，先把前两个数相加或先把后两个数相加，和不变。

\[
a + b + c = (a + b) + c = a + (b + c)
\]


例如： $36 + 48 + 52$ ，如果先算 $48 + 52 = 100$ ，再算 $36 + 100 = 136$ ，比按顺序算更方便。

\end{understand}

\begin{workedexamples}


\textbf{例 1.} 用简便方法计算 $67 + 54 + 33$ 。

\textbf{解：} 观察到 $67 + 33 = 100$ ，利用交换律和结合律：

\[
67 + 54 + 33 = (67 + 33) + 54 = 100 + 54 = 154
\]


\textbf{例 2.} 用简便方法计算 $125 + 98$ 。

\textbf{解：} 把 $98$ 看作 $100 - 2$ ：

\[
125 + 98 = 125 + 100 - 2 = 225 - 2 = 223
\]


\end{workedexamples}

\begin{keytakeaway}


\begin{itemize}
  \item 加法交换律： $a + b = b + a$ 。
  \item 加法结合律： $(a + b) + c = a + (b + c)$ 。
  \item 利用这两条规律，可以调整加数的顺序和组合，使计算更简便。
\end{itemize}


\end{keytakeaway}

\begin{practice}


\begin{enumerate}
  \item 用简便方法计算： $88 + 46 + 12$ 。
  \item 用简便方法计算： $276 + 199$ 。
  \item 用简便方法计算： $75 + 168 + 25 + 32$ 。
\end{enumerate}


\end{practice}

\section*{乘法交换律与结合律}


\begin{explore}


体育课上，同学们站成 $4$ 行 $8$ 列做操。一共有多少人？不管你先数行还是先数列，结果一样吗？

\end{explore}

\begin{understand}


\textbf{乘法交换律}：两个数相乘，交换因数的位置，积不变。

\[
a \times b = b \times a
\]


\textbf{乘法结合律}：三个数相乘，先把前两个数相乘或先把后两个数相乘，积不变。

\[
a \times b \times c = (a \times b) \times c = a \times (b \times c)
\]


例如： $25 \times 7 \times 4$ ，先算 $25 \times 4 = 100$ ，再算 $100 \times 7 = 700$ ，比按顺序算简单得多。

\end{understand}

\begin{workedexamples}


\textbf{例 1.} 用简便方法计算 $125 \times 32$ 。

\textbf{解：} 把 $32$ 拆成 $8 \times 4$ ：

\[
125 \times 32 = 125 \times 8 \times 4 = 1000 \times 4 = 4000
\]


\textbf{例 2.} 用简便方法计算 $25 \times 16$ 。

\textbf{解：} 把 $16$ 拆成 $4 \times 4$ ：

\[
25 \times 16 = 25 \times 4 \times 4 = 100 \times 4 = 400
\]


\end{workedexamples}

\begin{keytakeaway}


\begin{itemize}
  \item 乘法交换律： $a \times b = b \times a$ 。
  \item 乘法结合律： $(a \times b) \times c = a \times (b \times c)$ 。
  \item 常见的"好搭档"： $25 \times 4 = 100$ ， $125 \times 8 = 1000$ 。
\end{itemize}


\end{keytakeaway}

\begin{practice}


\begin{enumerate}
  \item 用简便方法计算： $8 \times 47 \times 125$ 。
  \item 用简便方法计算： $25 \times 24$ 。
  \item 用简便方法计算： $5 \times 13 \times 20$ 。
\end{enumerate}


\end{practice}

\section*{乘法分配律}


\begin{explore}


学校买了 $30$ 套桌椅，每张桌子 $45$ 元，每把椅子 $25$ 元。一共花了多少钱？你可以先算每套桌椅多少钱再乘以 $30$ ，也可以分别算桌子和椅子各多少钱再加起来。两种算法结果一样吗？

\end{explore}

\begin{understand}


\textbf{乘法分配律}：一个数与两个数的和相乘，等于这个数分别与这两个数相乘再相加。

\[
a \times (b + c) = a \times b + a \times c
\]


反过来也成立：

\[
a \times b + a \times c = a \times (b + c)
\]


这条规律也适用于减法：

\[
a \times (b - c) = a \times b - a \times c
\]


乘法分配律是最常用的简便运算工具。

\end{understand}

\begin{workedexamples}


\textbf{例 1.} 用简便方法计算 $102 \times 45$ 。

\textbf{解：} 把 $102$ 拆成 $100 + 2$ ：

\[
\begin{aligned}
102 \times 45 &= (100 + 2) \times 45 \\
&= 100 \times 45 + 2 \times 45 \\
&= 4500 + 90 \\
&= 4590
\end{aligned}
\]


\textbf{例 2.} 用简便方法计算 $99 \times 37$ 。

\textbf{解：} 把 $99$ 看作 $100 - 1$ ：

\[
\begin{aligned}
99 \times 37 &= (100 - 1) \times 37 \\
&= 100 \times 37 - 1 \times 37 \\
&= 3700 - 37 \\
&= 3663
\end{aligned}
\]


\end{workedexamples}

\begin{keytakeaway}


\begin{itemize}
  \item 乘法分配律： $a \times (b + c) = a \times b + a \times c$ 。
  \item 反向使用（提取公因数）同样重要： $a \times b + a \times c = a \times (b + c)$ 。
  \item 遇到接近整十、整百的数，把它拆成"整数 $\pm$ 小数"再分配。
\end{itemize}


\end{keytakeaway}

\begin{practice}


\begin{enumerate}
  \item 用简便方法计算： $35 \times 98$ 。
  \item 用简便方法计算： $72 \times 56 + 72 \times 44$ 。
  \item 用简便方法计算： $101 \times 67$ 。
  \item 用简便方法计算： $48 \times 25$ （提示： $48 = 50 - 2$ 或 $25 \times 4 = 100$ ）。
\end{enumerate}


\end{practice}



\begin{answer}


\textbf{练一练 加法交换律与结合律}

\begin{enumerate}
  \item $88 + 46 + 12 = (88 + 12) + 46 = 100 + 46 = 146$ 。
  \item $276 + 199 = 276 + 200 - 1 = 476 - 1 = 475$ 。
  \item $75 + 168 + 25 + 32 = (75 + 25) + (168 + 32) = 100 + 200 = 300$ 。
\end{enumerate}


\textbf{练一练 乘法交换律与结合律}

\begin{enumerate}
  \item $8 \times 47 \times 125 = (8 \times 125) \times 47 = 1000 \times 47 = 47000$ 。
  \item $25 \times 24 = 25 \times 4 \times 6 = 100 \times 6 = 600$ 。
  \item $5 \times 13 \times 20 = (5 \times 20) \times 13 = 100 \times 13 = 1300$ 。
\end{enumerate}


\textbf{练一练 乘法分配律}

\begin{enumerate}
  \item $35 \times 98 = 35 \times (100 - 2) = 3500 - 70 = 3430$ 。
  \item $72 \times 56 + 72 \times 44 = 72 \times (56 + 44) = 72 \times 100 = 7200$ 。
  \item $101 \times 67 = (100 + 1) \times 67 = 6700 + 67 = 6767$ 。
  \item 方法一： $48 \times 25 = (50 - 2) \times 25 = 1250 - 50 = 1200$ 。方法二： $48 \times 25 = 12 \times (4 \times 25) = 12 \times 100 = 1200$ 。
\end{enumerate}


\end{answer}
